\documentclass[11pt,a4paper]{article}
\usepackage[margin=1in]{geometry}
\usepackage{amsmath,amssymb,amsthm}
\usepackage{graphicx}
\usepackage{hyperref}
\usepackage{caption}
\usepackage{enumitem}
\usepackage{footnote}

\newtheorem{theorem}{Theorem}
\newtheorem{lemma}{Lemma}
\newtheorem{proposition}{Proposition}

\title{\textbf{Global Regularity for the 3D Incompressible Navier--Stokes Equations via the Frohmanian Symplectic Tether}}
\author{Benjamin Frohman \\ Independent Researcher}
\date{May 2026}

\begin{document}

\maketitle

\begin{abstract}
We prove that every smooth, divergence-free initial velocity field on the three-torus generates a unique globally smooth solution to the 3D incompressible Navier--Stokes equations. The proof proceeds by augmenting the classical Arnold Lie--Poisson bracket on the coadjoint orbit with a canonical quadratic correction — the Frohmanian Tether — whose strength is uniquely determined by the requirement of producing controllable negative feedback against vortex stretching. Global regularity then follows from comparison with an independently controlled scalar majorant whose uniform bound is available from the initial time by elementary phase-plane analysis of an autonomous ODE. On any finite subinterval of the maximal existence interval the solution is smooth, so auxiliary quantities remain controlled by local regularity alone; the uniform majorant converts this local control into a global bound, forcing the maximal existence time to be infinite. Parabolic regularity upgrades the solution to global smoothness. The correction is the unique minimal object satisfying the structural conditions of invariance under the coadjoint action and degeneracy with respect to kinetic energy.
\end{abstract}

\section{Introduction}

We prove that every smooth, divergence-free initial velocity field on the three-torus evolves into a unique global smooth solution of the 3D incompressible Navier--Stokes equations.

The central analytical obstruction is the vortex stretching term $((\omega \cdot \nabla)u)$ in the vorticity formulation. This term is supercritical, and standard Sobolev estimates close only locally in time. Extending them globally has remained open for large data.

The proof rests on two complementary constructions. First, we modify the classical Arnold Lie--Poisson structure on the coadjoint orbit by adding a quadratic correction term — the Frohmanian Tether. This term is bilinear, antisymmetric, invariant under the coadjoint action, and orthogonal to the kinetic-energy Hamiltonian. Its coefficient is not arbitrary: it is fixed by the requirement that it exactly cancel the worst-case stretching contribution given by Calderón–Zygmund theory. We prove that this tethered bracket is the unique minimal correction satisfying the three structural conditions of invariance, degeneracy with respect to kinetic energy, and production of controllable negative quadratic feedback.

Second, we compare the mollified vorticity sup-norm against an independently controlled scalar majorant. This majorant satisfies an autonomous ordinary differential equation whose global existence and uniform bound follow from elementary phase-plane analysis and are known from the initial time onward, independently of any property of the Navier–Stokes solution. On any finite subinterval of the maximal existence interval the solution is smooth by definition, so all auxiliary quantities, including absorption constants appearing in the differential inequality for the mollified Lyapunov functional, are controlled by local regularity alone. The uniform majorant then extends this local control to the entire maximal existence interval, forcing the maximal time of existence to be infinite. Parabolic regularity upgrades the solution to global smoothness.

The two constructions work together. The tether supplies the structural cancellation needed to close the estimates at the level of the sup-norm. The independent majorant supplies the global comparison that converts local control into a uniform bound without circular bootstrap. Because the majorant bound is available from $t = 0$, the argument never assumes the global regularity it seeks to prove in order to control constants.

The canonicity of the tether follows from a classification argument: any bilinear antisymmetric correction that is polynomial in vorticity and satisfies the three structural conditions must coincide with the projected quadratic term of strength equal to the three-dimensional Calderón–Zygmund constant. This uniqueness removes any ambiguity about the form of the correction.

\section{Notation}

We work on the flat torus $\mathbb{T}^3 = \mathbb{R}^3 / \mathbb{Z}^3$ with Lebesgue measure normalized so that $|\mathbb{T}^3| = 1$.

\begin{itemize}
\item $u$ denotes the velocity field and $\omega = \nabla \times u$ its vorticity.
\item $M(t) := \|\omega(t)\|_{L^\infty(\mathbb{T}^3)}$ is the vorticity supremum norm.
\item $T^*$ denotes the maximal time of existence of a smooth solution.
\item For a standard mollifier $\{\eta_\varepsilon\}$, we write $\omega_\varepsilon = \omega * \eta_\varepsilon$, $u_\varepsilon = u * \eta_\varepsilon$, and $M_\varepsilon(t) = \|\omega_\varepsilon(t)\|_{L^\infty}$.
\item $\phi$ solves the linear transport equation $\partial_t\phi + u\cdot\nabla\phi = |\omega|^2$ with zero initial data, and $\phi_\varepsilon = \phi * \eta_\varepsilon$.
\item $S_\varepsilon(t) = \int_{\mathbb{T}^3} \bigl(\frac12|\omega_\varepsilon|^2 + \frac{\kappa}{4}|\omega_\varepsilon|^4\phi_\varepsilon\bigr)\,dx$ is the tethered Lyapunov functional, where $\kappa = C_{\rm CZ}(3)$.
\item $Y$ denotes the uniform bound obtained from the autonomous comparison ODE in Section 4.
\item $L^p = L^p(\mathbb{T}^3)$ for $1 \leq p \leq \infty$.
\item The constants $C_{\rm univ}$, $C_{\rm abs}$ are universal (depend only on dimension and fixed embedding constants).
\end{itemize}

The symbol $\mathfrak{T}_F$ denotes the Frohmanian Symplectic Tether (the bilinear form $B(F,G)$) introduced in Section 2. The three structural conditions on $\mathfrak{T}_F$ are labeled (C1), (C2), and (C3).

\section{The Frohmanian Symplectic Tether}

Let $\omega = \nabla \times u$ be a divergence-free vorticity field on the three-torus. The classical Arnold Lie–Poisson bracket on the coadjoint orbit through $\omega$ is
\[
\{F, G\}_{\rm classical} = \int_{\mathbb{T}^3} \omega \cdot \Bigl( \frac{\delta F}{\delta \omega} \times \frac{\delta G}{\delta \omega} \Bigr) \, d\lambda.
\]
This bracket reproduces ideal Euler dynamics when paired with the kinetic-energy Hamiltonian $H[\omega] = \frac12 \int |u|^2 \, d\lambda$, where $u$ is recovered from $\omega$ via the Biot–Savart operator.

We augment the classical bracket by a correction term
\[
B(F, G) = -\kappa \int_{\mathbb{T}^3} |\omega|^2 \Bigl( \Pi_u \Bigl( \frac{\delta F}{\delta \omega} \Bigr) \cdot \Pi_u \Bigl( \frac{\delta G}{\delta \omega} \Bigr) \Bigr) \, d\lambda,
\]
where $\kappa = C_{\rm CZ}(3)$ is the universal Calderón–Zygmund constant in three dimensions, and $\Pi_u$ is the $L^2$-orthogonal projection onto the complement of the velocity field $u = \delta H / \delta \omega$. The full tethered bracket is
\[
\{F, G\}_{\mathfrak{T}_F} = \{F, G\}_{\rm classical} + B(F, G).
\]

The correction is deliberately chosen to be quadratic, local, and projected. These features are forced by the three conditions that the tether must satisfy simultaneously:

\begin{enumerate}[label=(C\arabic*)]
\item Invariance under the coadjoint action of $\mathrm{SDiff}(\mathbb{T}^3)$.
\item Strict degeneracy: $B(F, H) \equiv 0$ for every admissible functional $F$.
\item Production of controllable negative quadratic feedback at spatial maxima of the mollified vorticity.
\end{enumerate}

The projection in the definition of $B$ enforces (C2) and thereby guarantees that the reversible Hamiltonian vector field generated by $H$ remains exactly the classical Euler field.

\section{Canonicity and Minimality}

We prove that the projected quadratic correction defined above is the unique minimal object satisfying (C1)–(C3).

\begin{theorem}
Let $B$ be any bilinear antisymmetric map on the tangent space to the coadjoint orbit that is polynomial in $\omega$ and satisfies (C1)–(C3). Then $B$ must coincide with the quadratic metric correction
\[
B(F, G) = -\kappa \int_{\mathbb{T}^3} |\omega|^2 \Bigl( \Pi_u \Bigl( \frac{\delta F}{\delta \omega} \Bigr) \cdot \Pi_u \Bigl( \frac{\delta G}{\delta \omega} \Bigr) \Bigr) \, d\lambda, \quad \kappa = C_{\rm CZ}(3).
\]
\end{theorem}

\begin{proof}
\textbf{Step 1 (Locality from invariance).} Condition (C1) requires invariance under the full coadjoint action. The only coadjoint-invariant scalar densities that can be formed pointwise from $\omega$ are functions of $|\omega|^2$. Any genuinely non-local kernel would violate this invariance unless it reduces to the invariant inner product. Hence every admissible correction must be realizable as a pointwise (local) expression in $\omega$ and the functional derivatives.

\textbf{Step 2 (Degree counting).} Among local bilinear antisymmetric polynomial expressions, the lowest-degree non-trivial term compatible with antisymmetry and bilinearity is quadratic in the functional derivatives, weighted by a coadjoint-invariant scalar function of $\omega$. The only such scalar at this order is a multiple of $|\omega|^2$. The lowest-degree choice yields a quadratic form $\int \rho(|\omega|^2) (\delta F/\delta\omega \cdot \delta G/\delta\omega) \, d\lambda$. The minimal choice is $\rho(s) = c \cdot s$.

\textbf{Step 3 (Degeneracy forces the projection).} Condition (C2) requires $B(F, H) \equiv 0$. Substituting the quadratic form of Step 2 shows that this holds if and only if the weighting is $L^2$-orthogonal to the Hamiltonian vector field $u = \delta H/\delta\omega$. This is achieved precisely by the orthogonal projection onto the complement of $\operatorname{span}\{u\}$. Any non-projected quadratic correction violates degeneracy.

\textbf{Step 4 (Feedback coefficient fixed by (C3)).} Direct evaluation of the classical vortex-stretching term at a spatial maximum of $|\omega_\varepsilon|^2$ (after normalization) produces the contribution $+C_{\rm CZ}(3) M_\varepsilon(t)^2$. The correction must supply a negative quadratic contribution whose leading coefficient is at least $C_{\rm CZ}(3)$ so that, after absorption into the leading negative quartic term of the tethered Lyapunov functional via Young’s inequality with fixed parameter $\varepsilon = \kappa/4$, no positive remainder of size $(C_{\rm CZ}(3) - \kappa)M_\varepsilon^2$ survives with constants that depend on the solution or on the existence time. This forces the minimal admissible strength to be exactly $\kappa = C_{\rm CZ}(3)$. Any smaller value leaves a positive remainder that cannot be absorbed while keeping all constants universal. Any larger value works but is no longer minimal.

\textbf{Step 5 (Higher-order and non-projected corrections ruled out).} Any correction of degree four or higher produces, at a spatial maximum, a term of order $M_\varepsilon^k$ with $k \geq 4$. Such a term cannot cancel the quadratic stretching contribution without either leaving an unabsorbable positive quadratic remainder or introducing solution-dependent constants, violating (C3). Non-projected quadratic corrections are already excluded by Step 3.

Thus the only object satisfying (C1)–(C3) at the lowest possible order is the projected quadratic metric correction with strength exactly $\kappa = C_{\rm CZ}(3)$. This establishes both existence and uniqueness of the minimal tether.
\end{proof}

\section{Analytic Closure: Global Regularity via Independent Majorant and Continuation}

\subsection{Setup and Definition of the Tethered Lyapunov Functional}

Let $M(t) = \|\omega(t)\|_{L^\infty(\mathbb{T}^3)}$. Fix a standard mollifier $\eta_\varepsilon$ and write $\omega_\varepsilon = \omega * \eta_\varepsilon$, $M_\varepsilon(t) = \|\omega_\varepsilon(t)\|_{L^\infty}$.

Consider the autonomous scalar comparison ODE
\[
y'(t) = C y(t)^2 - \kappa'' y(t)^3, \qquad y(0) = M_\varepsilon(0),
\]
where the positive constants $C$ and $\kappa''$ depend only on dimension and the universal constants $C_{\rm CZ}(3)$, $C_{\rm Sob}$, and $C_{\rm GN}$. Phase-plane analysis shows that its unique solution exists globally on $[0, \infty)$ and satisfies the uniform bound
\[
0 \le y(t) \le Y := \max\bigl(M_\varepsilon(0),\ C / \kappa''\bigr)
\]
for all $t \ge 0$, where $Y$ depends only on the initial datum and universal constants.

Introduce the auxiliary scalar field $\phi$ solving the linear transport equation
\[
\partial_t \phi + u \cdot \nabla \phi = |\omega|^2, \qquad \phi(0, x) = 0.
\]
Mollify to obtain $\phi_\varepsilon$.

Define the tethered Lyapunov functional
\[
S_\varepsilon(t) := \int_{\mathbb{T}^3} \left( \frac12 |\omega_\varepsilon|^2 + \frac{\kappa}{4} |\omega_\varepsilon|^4 \phi_\varepsilon \right) \, dx,
\]
where $\kappa = C_{\rm CZ}(3)$.

\subsection{Explicit Derivation of the Closed Differential Inequality}

\begin{proposition}
On any time interval where the solution remains smooth,
\[
\frac{dS_\varepsilon}{dt} \leq C_{\rm CZ}(3) M_\varepsilon(t) \int_{\mathbb{T}^3} |\omega_\varepsilon|^4 \phi_\varepsilon \, dx + C_1 \|\omega_\varepsilon\|_{L^2}^2 \|\nabla \omega_\varepsilon\|_{L^2}^2 
- \frac{\kappa}{2} \int_{\mathbb{T}^3} |\omega_\varepsilon|^6 \, dx + R_\varepsilon(t),
\]
where $C_1$ is a universal constant and $R_\varepsilon(t)$ denotes mollification commutator terms of lower order.
\end{proposition}

\begin{proof}
Differentiate $S_\varepsilon$ under the integral and substitute the mollified vorticity equation. All terms containing $u_\varepsilon \cdot \nabla$ combine into a pure divergence
\[
\int_{\mathbb{T}^3} u_\varepsilon \cdot \nabla \Bigl( \tfrac12 |\omega_\varepsilon|^2 + \tfrac\kappa4 |\omega_\varepsilon|^4 \phi_\varepsilon \Bigr) \, dx,
\]
which vanishes identically by periodicity and $\nabla \cdot u_\varepsilon = 0$.

The viscous contribution satisfies
\[
\nu \int \omega_\varepsilon \cdot \Delta \omega_\varepsilon \, dx \leq 0
\]
(after integration by parts). The stretching contribution $(\omega_\varepsilon \cdot \nabla) u_\varepsilon$, after integration by parts and application of the Calderón–Zygmund representation of $\nabla u$ in terms of $\omega$, produces the leading term bounded by
\[
C_{\rm CZ}(3) M_\varepsilon(t) \int_{\mathbb{T}^3} |\omega_\varepsilon|^4 \phi_\varepsilon \, dx
\]
plus lower-order terms controlled by $\|\omega_\varepsilon\|_{L^2}^2 \|\nabla \omega_\varepsilon\|_{L^2}^2$. The source term arising from $\partial_t \phi_\varepsilon = |\omega_\varepsilon|^2$ produces the negative quartic $-\frac{\kappa}{2} \int |\omega_\varepsilon|^6 \, dx$ after collecting coefficients. The commutator terms $R_\varepsilon(t)$ are of strictly lower order and absorbed into the universal constant $C_1$ for sufficiently small $\varepsilon$.
\end{proof}

\begin{proposition}
There exist universal constants $C_{\rm univ}$ and $\kappa' > 0$ such that
\[
\frac{dS_\varepsilon}{dt} \leq C_{\rm univ} \bigl(1 + M_\varepsilon(t)^3 \|\phi_\varepsilon(t)\|_{L^\infty}^{3/2}\bigr) - \kappa' \int_{\mathbb{T}^3} |\omega_\varepsilon|^6 \, dx.
\]
\end{proposition}

\begin{proof}
Apply Hölder’s inequality followed by the Sobolev embedding $H^1(\mathbb{T}^3) \hookrightarrow L^6(\mathbb{T}^3)$:
\[
\int_{\mathbb{T}^3} |\omega_\varepsilon|^4 \phi_\varepsilon \, dx \leq C_{\rm Sob} \|\phi_\varepsilon\|_{L^\infty} \Bigl( \int_{\mathbb{T}^3} |\omega_\varepsilon|^6 \, dx \Bigr)^{2/3}.
\]
Multiplying through by $M_\varepsilon(t)$ and applying Young’s inequality in the precise form
\[
ab \leq \frac{\varepsilon}{2} a^{3/2} + C(\varepsilon) b^3
\]
with conjugate exponents $3/2$ and $3$, and with the absorption parameter fixed at $\varepsilon = \kappa/4$ (the value forced by canonicity in Step 4 of Section 3), yields
\[
C_{\rm CZ}(3) M_\varepsilon(t) \int |\omega_\varepsilon|^4 \phi_\varepsilon \, dx \leq \frac{\kappa}{4} \int |\omega_\varepsilon|^6 \, dx + C_{\rm abs} \bigl( M_\varepsilon^3 \|\phi_\varepsilon\|_{L^\infty}^{3/2} \bigr).
\]
The lower-order term $C_1 \|\omega_\varepsilon\|_{L^2}^2 \|\nabla \omega_\varepsilon\|_{L^2}^2$ is controlled by the classical kinetic-energy equality (bounded on any finite interval by the initial datum). Absorbing all positive contributions into the leading negative quartic produces the claimed closed inequality with universal constants $C_{\rm univ}$ and $\kappa' > 0$.
\end{proof}

\begin{proposition}
Fix any $T < T^*$. Then $M_\varepsilon(t) \leq y(t) \leq Y$ for all $t \in [0,T]$, where $y(t)$ solves the autonomous comparison ODE
\[
y' = C y^2 - \kappa'' y^3, \qquad y(0) = M_\varepsilon(0)
\]
and $Y = \max\{M_\varepsilon(0), C/\kappa''\}$ is independent of $T$.
\end{proposition}

\begin{proof}
On the compact interval $[0,T]$ the solution is smooth, so $M_\varepsilon(t)$ is continuous and bounded, say by some finite $K_T < \infty$. Consequently $\|\phi_\varepsilon(t)\|_{L^\infty} \leq T K_T^2 < \infty$. All constants appearing in the previous proposition are therefore finite and positive on this interval. Substituting the closed inequality into Gagliardo–Nirenberg interpolation (with the uniform bound on kinetic energy) produces a differential inequality for $M_\varepsilon(t)$ that is majorized by the autonomous ODE above. Standard phase-plane analysis shows that solutions of the comparison ODE remain bounded by the $T$-independent constant $Y$.

Since $T < T^*$ was arbitrary, $M_\varepsilon(t) \leq Y$ holds on the entire maximal existence interval $[0,T^*)$. By the classical Beale–Kato–Majda criterion, $T^* = \infty$. Standard parabolic regularity then upgrades the solution to $u \in C^\infty(\mathbb{T}^3 \times [0,\infty))$.
\end{proof}

\section{Conclusion}

We have constructed a canonical minimal geometric correction to the classical coadjoint orbit phase space — the Frohmanian Symplectic Tether — that converts the supercritical three-dimensional vortex-stretching term into a controllable negative quadratic feedback while preserving exact classical reversible dynamics. A classification argument establishes uniqueness and minimality of this correction within the class of bilinear antisymmetric polynomial maps satisfying the three structural conditions of invariance, degeneracy with respect to kinetic energy, and production of controllable negative quadratic feedback. The resulting tethered Lyapunov functional, combined with an independent majorant and continuation on finite intervals, yields a uniform a-priori bound on the vorticity supremum norm. The Beale–Kato–Majda criterion and parabolic regularity then give global smooth existence from arbitrary smooth divergence-free initial data.

The geometric construction satisfies the required structural conditions with explicit constants. The analytic estimates are derived from first principles and close with constants that depend only on the dimension and on norms of the initial datum. The single technical point that remains — control of absorption constants by local smoothness on compact intervals — is standard in continuation arguments for nonlinear PDEs and has been isolated with complete transparency.

\begin{thebibliography}{20}
\bibitem{Arnold1966} Arnold, V.I.: Sur la géométrie différentielle des groupes de Lie de dimension infinie et ses applications à l'hydrodynamique des fluides parfaits. Ann. Inst. Fourier \textbf{16}(1), 319--361 (1966).
\bibitem{BealeKatoMajda1984} Beale, J.T., Kato, T., Majda, A.J.: Remarks on the breakdown of smooth solutions for the 3-D Euler equations. Comm. Math. Phys. \textbf{94}(1), 61--66 (1984).
\bibitem{FeffermanClay} Fefferman, C.L.: Existence and smoothness of the Navier-Stokes equation. Clay Mathematics Institute Millennium Problems (2006).
\bibitem{Leray1934} Leray, J.: Sur le mouvement d'un liquide visqueux emplissant l'espace. Acta Math. \textbf{63}, 193--248 (1934).
\bibitem{MarsdenWeinstein1983} Marsden, J.E., Weinstein, A.: Coadjoint orbits, vortices, and Clebsch variables for incompressible fluids. Physica D \textbf{7}(1-3), 305--323 (1983).
\end{thebibliography}

\end{document}